% =============================================================
% 中英文摘要（中文期刊惯例：双语并存）
% =============================================================

\begin{abstract}
\textbf{中文摘要。}大语言模型（LLM）推理负载的快速增长带来了服务成本与能耗压力，单轴吞吐量导向的调度器已无法应对。本文研究 AI 生成内容（AIGC）服务特有的五项物理约束下的\emph{云--边 LLM 推理调度}：模型权重驻留、键值（KV）缓存本地性、连续批处理、解码阶段的显存带宽下限，以及两阶段请求生命周期。我们设计了一个 AIGC 感知的近端策略优化（PPO）调度器，通过状态特征与奖励整形将上述物理量显式暴露给策略网络，并与涵盖启发式、负载均衡、bandit、多目标搜索与学习式方法的 10 个基线对比。在覆盖拓扑、负载与工作负载组合的 30 种配置（每种 $N{=}30$ 次运行）下，本文调度器于典型工况在均值性能上取得严格 Pareto 占优：相对最强单目标基线，服务等级目标（\slo{}）达成率提升 28\%（统计显著，$p<0.05$），同时每 token 能耗降低 7\%。即便面对总体最强基线---膝点 NSGA-II 搜索---本文调度器在两项均值指标上仍更优，且在中等负载下差异显著（$p<0.01$），在每一种配置下每 token 能耗都更低，而调度时计算开销少 $5\times$。然而，12 组件消融揭示联合增益是一项系统级整体贡献：仅连续批处理仿真与充足预训练个体显著；所有 AIGC 奖励与状态组件在单独消融时均与完整模型在统计上不可区分。我们开源仿真器及 30+ 套基准 manifest，以支持后续 AIGC 调度研究。

\textbf{关键词：}LLM 推理；AIGC；云边协同；强化学习；调度；能效；Pareto 优化。

\textbf{中图分类号：}TP18

\textbf{文献标识码：}A
\end{abstract}

% --- English Abstract（中文期刊要求双语并存） -----------------
\noindent\textbf{Abstract.} The rapid growth of large language
model (LLM) inference workloads
imposes serving-cost and energy pressures that single-axis
throughput-oriented schedulers cannot address. We study
\emph{cloud-edge LLM inference scheduling} under five
physical constraints specific to AI-generated-content (AIGC)
serving: model-weight residency,
key-value (KV) cache locality, continuous batching, the
memory-bandwidth floor of decode, and the two-phase
request lifecycle. We design an AIGC-aware
proximal policy optimization (PPO) scheduler that exposes these
physics to the policy through state features and reward
shaping, and compare it against ten baselines spanning
heuristic, load-balancing, bandit, multi-objective search,
and learned approaches.
Across 30 configurations sweeping topology, load, and workload
composition ($N{=}30$ runs each), our scheduler achieves
strict Pareto dominance on mean performance in the canonical
operating regime: a
28\% gain in service-level-objective (\slo) attainment
(statistically significant,
$p<0.05$) together with 7\% lower energy per token relative to
the strongest single-objective baseline. Even against a
knee-point NSGA-II search---the strongest baseline
overall---our scheduler is better on both means,
significantly so at moderate load ($p<0.01$), with lower
energy per token in every configuration at $5\times$ less
scheduling-time compute. A 12-component ablation,
however, reveals that the joint gain is a holistic system
contribution: only continuous-batching simulation and adequate
pretraining are individually significant; every AIGC reward and
state component is statistically indistinguishable from the full
model when ablated alone. We release the simulator and the 30+
benchmark manifests to support future AIGC scheduling research.

\noindent\textbf{Key words:} LLM inference; AIGC; cloud-edge
computing; reinforcement learning; scheduling; energy
efficiency; Pareto optimization.
